\ifdefined\gkver\else\def\gkver{student}\fi
\documentclass[\gkver]{gaokaozhenti}
\begin{document}

\chapter{2023年河北}

\begin{enumerate}
\item
%% number: 1
%% paperName: 2023年河北高考第 1 题 4 分
%% typeId: 1
%% type: 单选题
%% chapter: 原子和原子核
%% point: 玻尔模型
%% method: 无
%% score: 4
%% degree: 850
%% duplicateId: 0
%% body:
2022 年 8 月 30 日，国家航天局正式发布了“羲和号”太阳探测卫星国际上首次在轨获取的太阳 H$_{\alpha}$ 谱线精细结构。H$_{\alpha}$ 是氢原子巴耳末系中波长最长的谱线，其对应的能级跃迁过程为\xzanswer{D}

\onepicture[w=4.0cm]{figs/01.png}

\fourchoices[answer=D]
{从 $\infty$ 跃迁到 $n = 2$}
{从 $n = 5$ 跃迁到 $n = 2$}
{从 $n = 4$ 跃迁到 $n = 2$}
{从 $n = 3$ 跃迁到 $n = 2$}

%% answer: D
%% memo:
\memoanswer{
$H_{\alpha}$ 是氢原子巴耳末系中波长最长的谱线，根据 $\nu=\frac{c}{\lambda}$ 可知 $H_{\alpha}$ 是氢原子巴耳末系中频率最小的谱线，根据氢原子的能级图，利用玻尔理论中的频率条件 $h\nu=E_{n}-E_{2}$

可见能级差越小，频率越低，波长越长。故 $H_{\alpha}$ 对应的能级跃迁过程为从 $n=3$ 跃迁到 $n=2$。

故选 D。
}

\item
%% number: 2
%% paperName: 2023年河北高考第 2 题 4 分
%% typeId: 1
%% type: 单选题
%% chapter: 光学
%% point: 光的干涉
%% method: 无
%% score: 4
%% degree: 750
%% duplicateId: 0
%% body:
制造某型芯片所使用的银灰色硅片覆上一层厚度均匀的无色透明薄膜后，在自然光照射下硅片呈现深紫色。关于此现象，下列说法正确的是\xzanswer{C}

\fourchoices[answer=C]
{上述现象与彩虹的形成原理相同}
{光在薄膜的下表面发生了全反射}
{薄膜上下表面的反射光发生了干涉}
{薄膜厚度发生变化，硅片总呈现深紫色}

%% answer: C
%% memo:
\memoanswer{
A．上述现象是由于光的干涉造成的，彩虹的形成原理主要为光的折射，上述现象与彩虹的形成原理不相同，故 A 错误；

BC．硅片呈现深紫色的原因是薄膜的厚度正好使紫光在薄膜上下表面的反射光发生干涉，振动加强，故 B 错误，C 正确；

D．根据光的干涉中相互加强的条件，可知当薄膜的厚度发生变化时，满足加强条件的波长也会发生改变，导致硅片呈现不同的颜色，故 D 错误。

故选 C。
}

\item
%% number: 3
%% paperName: 2023年河北高考第 3 题 4 分
%% typeId: 1
%% type: 单选题
%% chapter: 曲线运动
%% point: 万有引力定律
%% method: 无
%% score: 4
%% degree: 550
%% duplicateId: 0
%% body:
我国自古就有“昼涨为潮，夜涨为汐”之说，潮汐是月球和太阳对海水的引力变化产生的周期性涨落现象，常用引潮力来解释。月球对海水的引潮力大小与月球质量成正比、与月地距离的 3 次方成反比，方向如图~\ref{2023河北03a}，随着地球自转，引潮力的变化导致了海水每天 2 次的潮涨潮落。太阳对海水的引潮力与月球类似，但大小约为月球引潮力的 0.45 倍。每月 2 次大潮（引潮力最大）和 2 次小潮（引潮力最小）是太阳与月球引潮力共同作用的结果，结合图~\ref{2023河北03b}，下列说法正确的是\xzanswer{A}

\twopicture[numA=1, numB=2, labelA=2023河北03a, labelB=2023河北03b, wA=6.5cm, wB=5.0cm]
{figs/03a.png}
{figs/03b.png}

\fourchoices[answer=A]
{月球在位置 1 时会出现大潮}
{月球在位置 2 时会出现大潮}
{涨潮总出现在白天，退潮总出现在夜晚}
{月球引潮力和太阳引潮力的合力一定大于月球引潮力}

%% answer: A
%% memo:
\memoanswer{
AB．太阳、月球、地球三者在同一条直线上，太阳和月球的引潮力叠加在一起，潮汐现象最明显，称为大潮，月地连线与日地连线互相垂直，太阳引潮力就会削弱月球的引潮力，形成小潮，如图~\ref{2023河北03b} 所示得月球在位置 1 时会出现大潮，故 A 正确，B 错误；

C．每一昼夜海水有两次上涨和两次退落，人们把每次在白天出现的海水上涨叫做“潮”，把夜晚出现的海水上涨叫做“汐”，合称潮汐，故 C 错误；

D．月地连线与日地连线互相垂直位置（2 位置）时月球引潮力和太阳引潮力的合力等于月球引潮力减太阳引潮力小于月球引潮力，故 D 错误。

故选 A。
}

\item
%% number: 4
%% paperName: 2023年河北高考第 4 题 4 分
%% typeId: 1
%% type: 单选题
%% chapter: 力物体的平衡
%% point: 三力平衡
%% method: 无
%% score: 4
%% degree: 650
%% duplicateId: 0
%% body:
如图，轻质细杆 AB 上穿有一个质量为 $m$ 的小球 C，将杆水平置于相互垂直的固定光滑斜面上，系统恰好处于平衡状态。已知左侧斜面与水平面成 $30^\circ$ 角，则左侧斜面对杆 AB 支持力的大小为\xzanswer{B}

\onepicture[w=5.5cm]{figs/04.png}

\fourchoices[answer=B]
{$mg$}
{$\frac{\sqrt{3}}{2}mg$}
{$\frac{\sqrt{3}}{3}mg$}
{$\frac{1}{2}mg$}

%% answer: B
%% memo:
\memoanswer{
对轻杆和小球组成的系统进行受力分析，如图设左侧斜面对杆 AB 支持力的大小为 $N_{A}$，由平衡条件有 $N_{A}=mg\cos 30^\circ$

得 $N_{A}=\frac{\sqrt{3}}{2}mg$

故选 B。

\onepicture[w=3.6cm]{figs/04_an.jpg}
}

\item
%% number: 5
%% paperName: 2023年河北高考第 5 题 4 分
%% typeId: 1
%% type: 单选题
%% chapter: 交流电
%% point: 变压器
%% method: 无
%% score: 4
%% degree: 550
%% duplicateId: 0
%% body:
除颤仪是用于突发心室纤颤等心脏疾病的急救医疗设备，某型除颤仪电路原理如图，某次调试时交流电压表示数为 $20\UV$，电容器充电完毕，开关由“1”掷向“2”，放电电流平均值为 $2.8\UA$，放电时间约为 $10^{-2}\Us$，已知电容器电容为 $2.0\times10^{-5}\UF$，则升压变压器原副线圈的匝数比约为\xzanswer{B}

\onepicture[w=8.0cm]{figs/05.png}

\fourchoices[answer=B]
{$1:25$}
{$1:50$}
{$1:70$}
{$1:100$}

%% answer: B
%% memo:
\memoanswer{
由于放电电流平均值为 $2.8\UA$，由电流的定义式有 $\overline{I}=\frac{Q}{t}$

根据电容器的电容定义有 $C=\frac{Q}{U}$

所以电容器两端的电压为 $U_{C}=1400\UV$

等于副线圈电压的最大值，所以副线圈两端的电压有效值为 $U_{2}=\frac{U_{C}}{\sqrt{2}}=700\sqrt{2}\UV$

由题意可知，升压变压器原线圈两端的电压有效值为 $U_{1}=20\UV$

由理想变压器有 $\frac{U_{1}}{U_{2}}=\frac{n_{1}}{n_{2}}$

解得 $\frac{n_{1}}{n_{2}}\approx\frac{1}{50}$

故选 B。
}

\item
%% number: 6
%% paperName: 2023年河北高考第 6 题 4 分
%% typeId: 1
%% type: 单选题
%% chapter: 运动和力
%% point: 动量定理
%% method: 无
%% score: 4
%% degree: 450
%% duplicateId: 0
%% body:
某科研团队通过传感器收集并分析运动数据，为跳高运动员的技术动作改进提供参考。图为跳高运动员在起跳过程中，其单位质量受到地面的竖直方向支持力随时间变化关系曲线。图像中 $10.10\Us$ 至 $10.35\Us$ 内，曲线下方的面积与阴影部分的面积相等。已知该运动员的质量为 $60\Ukg$，重力加速度 $g$ 取 $10\Umsq$。下列说法正确的是\xzanswer{C}

\onepicture[w=6.5cm]{figs/06.png}

\fourchoices[answer=C]
{起跳过程中运动员的最大加速度约为 $42\Umsq$}
{起跳后运动员重心上升的平均速度大小约为 $3\Ums$}
{起跳后运动员重心上升的最大高度约为 $0.45\Um$}
{起跳过程中运动员所受合力的冲量大小约为 $330\UNs$}

%% answer: C
%% memo:
\memoanswer{
A．由图像可知，运动员受到的最大支持力约为 $F_{\max}=42\times60\UN=2520\UN$

根据牛顿第二定律可知，起跳过程中运动员的最大加速度约为

$a_{\max}=\frac{F_{\max}-mg}{m}=\frac{2520-60\times10}{60}\Umsq=32\Umsq$

故 A 错误；

BCD．根据 $F-t$ 图像可知，起跳过程中支持力的冲量为

$I_{F}=22\times60\times(10.35-10.10)\UNs=330\UNs$

起跳过程中运动员所受合力的冲量大小约为

$I_{\text{合}}=I_{F}-mgt=330\UNs-60\times10\times(10.35-10.10)\UNs=180\UNs$

根据动量定理可得 $I_{\text{合}}=mv-0$

解得起跳离开地面瞬间的速度为 $v=3\Ums$

则起跳后运动员重心上升的平均速度为 $\overline{v}=\frac{v}{2}=1.5\Ums$

起跳后运动员重心上升的最大高度为

$h=\frac{v^{2}}{2g}=\frac{3^{2}}{2\times10}\Um=0.45\Um$

故 BD 错误，C 正确。

故选 C。
}

\item
%% number: 7
%% paperName: 2023年河北高考第 7 题 4 分
%% typeId: 1
%% type: 单选题
%% chapter: 运动和力
%% point: 动量和能量
%% method: 无
%% score: 4
%% degree: 450
%% duplicateId: 0
%% body:
由点电荷组成的系统的电势能与它们的电荷量、相对位置有关。如图~\ref{2023河北07a}，a、b、c，d 四个质量均为 $m$、带等量正电荷的小球，用长度相等、不可伸长的绝缘轻绳连接，静置在光滑绝缘水平面上，$O$ 点为正方形中心，设此时系统的电势能为 $E_{0}$。剪断 a、d 两小球间的轻绳后，某时刻小球 d 的速度大小为 $v$，方向如图~\ref{2023河北07b}，此时系统的电势能为\xzanswer{B}

\twopicture[numA=1, numB=2, labelA=2023河北07a, labelB=2023河北07b, wA=4.0cm, wB=6.5cm]
{figs/07a.png}
{figs/07b.png}

\fourchoices[answer=B]
{$E_{0}-\frac{7}{4}mv^{2}$}
{$E_{0}-\frac{5}{4}mv^{2}$}
{$E_{0}-\frac{1}{4}mv^{2}$}
{$E_{0}-\frac{5}{8}mv^{2}$}

%% answer: B
%% memo:
\memoanswer{
由于运动的对称性，所以 a 球的速度大小也为 $v$，方向与 ad 连线成 $30^\circ$ 角斜向左下方。b 和 c 球的运动方向垂直于 bc 向上，大小相等，将两球看成一个整体，其质量为 $2m$，速度为 $v_{1}$，对于四个球组成的系统来说，动量守恒，有

$0=2mv\sin 30^\circ-2mv_{1}$

解得 $v_{1}=\frac{1}{2}v$

由于整个系统的能量守恒，设此时系统的电势能为 $E_{1}$，有

$E_{0}=E_{1}+\frac{1}{2}\cdot2mv^{2}+\frac{1}{2}\cdot2mv_{1}^{2}$

解得 $E_{0}-\frac{5}{4}mv^{2}$

故选 B。
}

\item
%% number: 8
%% paperName: 2023年河北高考第 8 题 6 分
%% typeId: 2
%% type: 多选题
%% chapter: 电磁感应
%% point: 楞次定律推广
%% method: 无
%% score: 6
%% degree: 550
%% duplicateId: 0
%% body:
如图，绝缘水平面上四根完全相同的光滑金属杆围成矩形，彼此接触良好，匀强磁场方向竖直向下。金属杆 2、3 固定不动，1、4 同时沿图箭头方向移动，移动过程中金属杆所围成的矩形周长保持不变。当金属杆移动到图位置时，金属杆所围面积与初始时相同。在此过程中\xzanswer{CD}

\twopicture[showlabel=false, wA=5.6cm, wB=5.6cm]
{figs/08a.png}
{figs/08b.png}

\fourchoices[answer=CD]
{金属杆所围回路中电流方向保持不变}
{通过金属杆截面的电荷量随时间均匀增加}
{金属杆 1 所受安培力方向与运动方向先相同后相反}
{金属杆 4 所受安培力方向与运动方向先相反后相同}

%% answer: CD
%% memo:
\memoanswer{
A．由数学知识可知金属杆所围回路的面积先增大后减小，金属杆所围回路内磁通量先增大后减小，根据楞次定律可知电流方向先沿逆时针方向，后沿顺时针方向，故 A 错误；

B．由于金属杆所围回路的面积非均匀变化，故感应电流的大小不恒定，故通过金属杆截面的电荷量随时间不是均匀增加的，故 B 错误；

CD．由上述分析，再根据左手定则，可知金属杆 1 所受安培力方向与运动方向先相同后相反，金属杆 4 所受安培力方向与运动方向先相反后相同，故 CD 正确。

故选 CD。
}

\item
%% number: 9
%% paperName: 2023年河北高考第 9 题 6 分
%% typeId: 2
%% type: 多选题
%% chapter: 运动和力
%% point: 变加速直线运动
%% method: 无
%% score: 6
%% degree: 350
%% duplicateId: 0
%% body:
如图，质量为 $m$ 的小球穿在固定光滑杆上，与两个完全相同的轻质弹簧相连。开始时将小球控制在杆上的 $A$ 点，弹簧 1 竖直且处于原长，弹簧 2 处于水平伸长状态，两弹簧可绕各自转轴 $O_{1}$，$O_{2}$ 无摩擦转动。$B$ 为杆上的另一个点，与 $O_{1}$、$A$、$O_{2}$ 构成矩形，$AB=2AO_{1}$。现将小球从 $A$ 点释放，两弹簧始终处于弹性限度内。下列说法正确的是\xzanswer{BC}

\onepicture[w=6.5cm]{figs/09.png}

\fourchoices[answer=BC]
{小球沿杆在 $AB$ 之间做往复运动}
{与没有弹簧时相比，小球从 $A$ 点运动到 $B$ 点所用的时间更短}
{小球从 $A$ 点运动到 $B$ 点的过程中，两个弹簧对小球做的总功为零}
{小球从 $A$ 点运动到 $B$ 点的过程中，弹簧 2 的弹性势能先减小后增大}

%% answer: BC
%% memo:
\memoanswer{
AC．根据对称性可知，小球从 $A$ 点运动到 $B$ 点的过程中，两个弹簧对小球做的总功为零，则此过程合力做功等于重力对小球做的功，根据动能定理可知，小球在 $B$ 点的速度大于 0，所以小球到达 $B$ 点后继续向下运动，小球不会在 $AB$ 之间做简谐运动，故 A 错误，C 正确；

D．小球从 $A$ 点运动到 $B$ 点的过程中，弹簧 2 先从伸长状态变为原长，再从原长变为压缩状态，最后再恢复原长，故弹簧 2 的弹性势能先减小后增大再减小，故 D 错误；

B．小球从 $A$ 点运动到 $B$ 点过程，由于两个弹簧对小球做的总功为零，与没有弹簧时相比，小球运动到 $B$ 点的速度相等；没有弹簧时，小球运动的加速度为 $a=\frac{mg\sin\theta}{m}=g\sin\theta$。有弹簧时，加速度先大于 $g\sin\theta$，然后加速度逐渐减小，到 AB 中点时，加速度为 $g\sin\theta$，之后加速度小于 $g\sin\theta$，则两种情况的 $v-t$ 图像如图所示

两种情况的 $v-t$ 图像与横轴围成的面积相等，由图可知与没有弹簧时相比，小球从 $A$ 点运动到 $B$ 点所用的时间更短，故 B 正确。

故选 BC。

\onepicture[w=3.6cm]{figs/09_an.png}
}

\item
%% number: 10
%% paperName: 2023年河北高考第 10 题 6 分
%% typeId: 2
%% type: 多选题
%% chapter: 电场
%% point: 电势能电势差电场力做功
%% method: 无
%% score: 6
%% degree: 450
%% duplicateId: 0
%% body:
如图，在 $x$ 轴上放置四个点电荷 $q_{1}$、$q_{2}$、$q_{3}$ 和 $q_{4}$，$q_{1}$、$q_{2}$ 位于 A 点两侧，$q_{3}$、$q_{4}$ 位于 B 点两侧。C 点在 $y$ 轴上，且 OA = OB = OC。取无穷远处电势为零，电荷位置与电荷量满足一定关系，使得以 O 点为球心、以 OC 为半径的球面上各点的电势均为零。下列说法正确的是\xzanswer{BD}

\onepicture[w=7.0cm]{figs/10.png}

\fourchoices[answer=BD]
{A、B 两点电场强度的方向一定沿 $x$ 轴正方向}
{若在 O 点放一正点电荷，则 A、B 两点的电势一定升高}
{试探电荷 $q$ 沿 $y$ 轴运动过程中，若静电力始终不做功，则它受到的静电力始终为零}
{试探电荷 $q$ 沿 $y$ 轴运动过程中，若静电力始终不做功，则 C 点的电场强度一定为零}

%% answer: BD
%% memo:
\memoanswer{
A．以 O 点为球心、以 OC 为半径的球面为等势面，由电场强度和等势面的关系可知，A、B 两点电场强度的方向与半径垂直，与 $x$ 轴重合。由于无法判断各个电荷的电性，故 A、B 两点电场强度的方向无法判断，故 A 错误；

B．取无穷远处为零电势点，由于正点电荷周围的电势为正值，若在 O 点放一正点电荷，则 A、B 两点的电势一定升高，故 B 正确；

C．试探电荷 $q$ 沿 $y$ 轴运动过程中，根据电荷分布，若静电力始终不做功，则经过 $y$ 轴且垂直于 $x$ 轴的平面为等势面，但静电力不一定为零，故 C 错误；

D．根据以 O 点为球心、以 OC 为半径的球面为等势面，可知 C 点的电场强度方向应与 $y$ 轴重合，再根据经过 $y$ 轴且垂直于 $x$ 轴的平面为等势面，可知 C 点的电场强度方向应与轴垂直，同一点电场强度的方向是唯一的，故 C 点的电场强度一定为零，故 D 正确。

故选 BD。
}

\item
%% number: 11
%% paperName: 2023年河北高考第 11 题 8 分
%% typeId: 4
%% type: 实验题
%% chapter: 机械振动
%% point: 单摆测重力加速度
%% method: 无
%% score: 8
%% degree: 550
%% duplicateId: 0
%% body:
某实验小组利用图~\ref{2023河北11a}装置测量重力加速度。摆线上端固定在 O 点，下端悬挂一小钢球，通过光电门传感器采集摆动周期。
\begin{enumerate}
	\item
	关于本实验，下列说法正确的是\tkanswer{ABD}。（多选）

	A．小钢球摆动平面应与光电门 U 形平面垂直

	B．应在小钢球自然下垂时测量摆线长度

	C．小钢球可以换成较轻的橡胶球

	D．应无初速度、小摆角释放小钢球

	\item
	组装好装置，用毫米刻度尺测量摆线长度 $L$，用螺旋测微器测量小钢球直径 $d$。螺旋测微器示数如图~\ref{2023河北11b}，小钢球直径 $d=$ \tkanswer{20.035}$\Umm$，记摆长 $l=L+\frac{d}{2}$。

	\item
	多次改变摆线长度，在小摆角下测得不同摆长 $l$ 对应的小钢球摆动周期 $T$，并作出 $l-T^{2}$ 图像，如图~\ref{2023河北11c}。

	根据图线斜率可计算重力加速度 $g=$ \tkanswer{9.87}$\Umsq$（保留 3 位有效数字，$\pi^{2}$ 取 9.87）。

	\item
	若将摆线长度误认为摆长，仍用上述图像法处理数据，得到的重力加速度值将\tkanswer{不变}（填“偏大”“偏小”或“不变”）。
\end{enumerate}

\threepicture[labelA=2023河北11a, labelB=2023河北11b, labelC=2023河北11c, numA=1, numB=2, numC=3, showlabel=false, wA=4.0cm, wB=3.6cm, wC=5.2cm, align=right]
{figs/11a.png}
{figs/11b.png}
{figs/11c.png}

%% answer:
\jdanswer{
\begin{enumerate}
	\item
	ABD

	\item
	20.035

	\item
	9.87

	\item
	不变

\end{enumerate}
}
%% memo:
\memoanswer{
（1）A．使用光电门测量时，光电门 U 形平面与被测物体的运动方向垂直是光电门使用的基本要求，故 A 正确；

B．测量摆线长度时，要保证绳子处于伸直状态，故 B 正确；

C．单摆是一个理想化模型，若采用质量较轻的橡胶球，空气阻力对摆球运动的影响较大，故 C 错误；

D．无初速度、小摆角释放的目的是保持摆球在竖直平面内运动，不形成圆锥摆，且单摆只有在摆角很小的情况下才可视为简谐运动，使用 $T=2\pi\sqrt{\frac{l}{g}}$ 计算单摆的周期，故 D 正确。故选 ABD。

（2）小钢球直径为 $d=20\Umm+3.5\times0.01\Umm=20.035\Umm$

（3）单摆周期公式 $T=2\pi\sqrt{\frac{l}{g}}$

整理得 $l=\frac{g}{4\pi^{2}}T^{2}$

由图像知图线的斜率 $k=\frac{g}{4\pi^{2}}=\frac{1.10-0.60}{4.40-2.40}\Umsq=\frac{1}{4}\Umsq$

解得 $g=9.87\Umsq$

（4）若将摆线长度 $L$ 误认为摆长 $l$，有 $T=2\pi\sqrt{\frac{L+\frac{d}{2}}{g}}$

则得到的图线为 $L=\frac{gT^{2}}{4\pi^{2}}-\frac{d}{2}$

仍用上述图像法处理数据，图线斜率不变，仍为 $\frac{g}{4\pi^{2}}$，故得到的重力加速度值不变。
}

\item
%% number: 12
%% paperName: 2023年河北高考第 12 题 8 分
%% typeId: 4
%% type: 实验题
%% chapter: 电路
%% point: 测量电阻率
%% method: 无
%% score: 8
%% degree: 450
%% duplicateId: 0
%% body:
某实验小组测量绕制滑动变阻器电阻丝的电阻率 $\rho$，实验器材：电源 $E$（$3\UV$，内阻很小）、电压表（量程为 $0\sim3\UV$，内阻很大）、待测滑动变阻器 $R_{x}$（最大阻值几十欧姆）、电阻箱 $R_{0}$（最大阻值 $9999\UO$）、滑动变阻器 $R_{1}$（最大阻内值 $50\UO$）、毫米刻度尺、开关 S 以及导线若干。实验电路如图。

第一步：按图连好电路。将滑动变阻器 $R_{1}$ 和 $R_{x}$ 的滑片均置于最左端，电阻箱 $R_{0}$ 阻值调为零。闭合开关 S，调整 $R_{1}$ 滑片的位置，使电压表的示数为 $2.00\UV$。

第二步：断开开关 S，保持 $R_{1}$ 滑片的位置不动，将 $R_{0}$ 的阻值调为 $5\UO$。

第三步：闭合开关 S，向右移动 $R_{x}$ 的滑片 P，使电压表的示数仍为 $2.00\UV$，记录 $R_{0}$ 的阻值 $R$ 以及 $R_{x}$ 的滑片 P 到左端点 $a$ 之间的距离 $l$。

第四步：断开开关 S，保持 $R_{1}$ 滑片的位置不动，调节 $R_{0}$ 的阻值分别为 $10\UO$、$15\UO$、$\ldots$，重复第三步。

第五步：实验结束，整理仪器。

实验记录的部分数据见下表。

\begin{center}
\begin{tabular}{|c|c|c|c|c|c|}
\hline
组次 & 1 & 2 & 3 & 4 & 5 \\
\hline
$R/\UO$ & 0 & 5 & 10 & 15 & 20 \\
\hline
$l/\Umm$ & 0 & 20.3 & 36.8 & 60.5 & 81.1 \\
\hline
\end{tabular}
\end{center}

\begin{enumerate}
	\item
	上表中不合理的一组数据为\tkanswer{3}（填组次序号）。

	\item
	当 $l$ 为 $80.0\Umm$ 时，$R_{x}$ 的滑片 P 到 $a$ 之间电阻丝的匝数为 133，电阻丝的半径 $r=\tkanswer{0.30}\Umm$（保留 2 位有效数字）。

	\item
	用 $l=kR$ 拟合上表数据，得 $k$ 近似为 $4.0\times10^{-3}\Um/\UO$，测得单匝电阻丝周长为 $90.0\Umm$，则电阻丝的电阻率 $\rho=$ \tkanswer{$4.7\times10^{-7}$}$\UOm$（保留 2 位有效数字）

	\item
	若电阻率 $\rho$ 的测量值与参考值相比偏大，产生误差的原因可能是\tkanswer{B}。

	A．未考虑电压表内阻

	B．未考虑电阻丝绝缘层厚度

	C．未考虑电源内阻
\end{enumerate}

\onepicture[w=6.0cm, align=right]{figs/12.png}

%% answer:
\jdanswer{
\begin{enumerate}
	\item
	3

	\item
	0.30

	\item
	$4.7\times10^{-7}$

	\item
	B

\end{enumerate}
}
%% memo:
\memoanswer{
（1）根据实验原理可知，待测电阻接入电路中的电阻阻值与电阻箱阻值之和应保持不变，而根据电阻定律有 $R_{x}=\rho\frac{l_{x}}{S}$，即电阻箱阻值增加多少，相应的待测电阻的阻值应减小多少，根据电阻定律可知，待测电阻的减少量与其接入电路中长度的减小少量成正比，即在误差允许范围内，减少相同的量则长度的减少量也应基本相同，则滑片 P 到 $a$ 端长度的增加量应基本相同，根据表中数据可知，电阻箱每增加 $5\UO$，滑片 P 到 $a$ 端长度的增加量分别是 $20.3\Umm$、$16.5\Umm$、$23.7\Umm$、$20.6\Umm$，由此可知第 3 组数据不合理。

（2）当 $l$ 为 $80.0\Umm$ 时，$R_{x}$ 的滑片 P 到 $a$ 之间电阻丝的匝数为 133，则有 $l=2r\cdot n$

因此可得电阻丝的半径 $r=\frac{l}{2n}$

其中 $n=133$

代入解得 $r=0.30\Umm$

（3）用 $l=kR$ 拟合上表数据，得 $k$ 近似为 $4.0\times10^{-3}\Um/\UO$，而电阻箱的阻值等于待测电阻滑片 P 到 $a$ 端电阻阻值的大小，即

$R=\rho\frac{L}{S}=\rho\frac{L}{\pi r^{2}}$

而根据拟合公式有 $R=\frac{l}{k}$

而单匝电阻丝长度为 $l=2r$

单匝电阻丝周长 $L=90\Umm$

联立可得

$\rho=\frac{\pi r^{2}\cdot l}{k\cdot L}=\frac{2\times3.14\times(0.30\times10^{-3})^{3}}{4.0\times10^{-3}\times90\times10^{-3}}\UOm=4.7\times10^{-7}\UOm$

（4）A．由于待测电阻接入电路中的电阻与电阻箱串联，而二者之和始终保持不变，因此电压表内阻不会对该实验产生影响，故 A 错误；

B．根据该实验中电阻率的计算公式可知，若未考虑电阻丝绝缘层厚度，则会造成测量值偏大，故 B 正确；

C．由于该实验采用分压式接法，而与滑动变阻器 $R_{1}$ 右边所串联的并联部分的电路的总电阻始终不变因此回路中的总电阻始终不变，则电源内阻对该实验不产生影响，故 C 错误。故选 B。
}

\item
%% number: 13
%% paperName: 2023年河北高考第 13 题 8 分
%% typeId: 6
%% type: 计算题
%% chapter: 气体性质
%% point: 盖吕萨克定律
%% method: 无
%% score: 8
%% degree: 450
%% duplicateId: 0
%% body:
如图，某实验小组为测量一个葫芦的容积，在葫芦开口处竖直插入一根两端开口、内部横截面积为 $0.1\Ucmq$ 的均匀透明长塑料管，密封好接口，用氮气排空内部气体，并用一小段水柱封闭氮气。外界温度为 $300\UK$ 时，气柱长度 $l$ 为 $10\Ucm$；当外界温度缓慢升高到 $310\UK$ 时，气柱长度变为 $50\Ucm$。已知外界大气压恒为 $1.0\times10^{5}\UPa$，水柱长度不计。
\begin{enumerate}
	\item
	求温度变化过程中氮气对外界做的功；
	\item
	求葫芦的容积；
	\item
	试估算被封闭氮气分子的个数（保留 2 位有效数字）。已知 $1\Umol$ 氮气在 $1.0\times10^{5}\UPa$、$273\UK$ 状态下的体积约为 $22.4\UL$，阿伏伽德罗常数 $N_{A}$ 取 $6.0\times10^{23}\Umoln$。
\end{enumerate}

\onepicture[w=3.5cm, align=right]{figs/13.png}

%% answer:
\jdanswer{
\begin{enumerate}
	\item
	$0.4\UJ$

	\item
	$119\Ucmc$

	\item
	$2.9\times10^{21}$

\end{enumerate}
}
%% memo:
\memoanswer{
（1）由于水柱的长度不计，故封闭气体的压强始终等于大气压强。设大气压强为 $p_{0}$，塑料管的横截面积为 $S$，初、末态气柱的长度分别为 $l_{1}$、$l_{2}$，气体对外做的功为 $W$。根据功的定义有

$W=p_{0}S(l_{2}-l_{1})$

解得 $W=0.4\UJ$

（2）设葫芦的容积为 $V$，封闭气体的初、末态温度分别为 $T_{1}$、$T_{2}$，体积分别为 $V_{1}$、$V_{2}$，根据盖-吕萨克定律有

$\frac{V_{1}}{T_{1}}=\frac{V_{2}}{T_{2}}$

$V_{1}=V+Sl_{1}$

$V_{2}=V+Sl_{2}$

联立以上各式并代入题给数据得 $V=119\Ucmc$

（3）设在 $1.0\times10^{5}\UPa$、$273\UK$ 状态下，$1\Umol$ 氮气的体积为 $V_{0}$、温度为 $T_{0}$，封闭气体的体积为 $V_{3}$，被封闭氮气的分子个数为 $n$。根据盖-吕萨克定律有

$\frac{V+Sl_{1}}{T_{1}}=\frac{V_{3}}{T_{0}}$

其中

$n=\frac{V_{3}}{V_{0}}N_{A}$

联立以上各式并代入题给数据得 $n=2.9\times10^{21}$ 个
}

\item
%% number: 14
%% paperName: 2023年河北高考第 14 题 14 分
%% typeId: 6
%% type: 计算题
%% chapter: 磁场
%% point: 带电粒子在组合场中的运动
%% method: 无
%% score: 14
%% degree: 450
%% duplicateId: 0
%% body:
如图~\ref{2023河北14a}，真空中两金属板 $M$、$N$ 平行放置，板间电压 $U$ 连续可调，在金属板 $N$ 的中心 $C$ 处开有一小孔。$F$、$G$、$H$ 为正三角形 $CDE$ 各边中点，$DE$ 与金属板平行。三角形 $FGH$ 区域内存在匀强磁场，磁感应强度大小为 $B$，方向垂直纸面向里。一带电粒子从紧邻 $M$ 板的 $P$ 点由静止释放沿 $CG$ 方向进入磁场，一段时间后沿 $\angle CED$ 的角平分线方向从 $E$ 点离开。已知正三角形 $CDE$ 的边长为 $a$，粒子质量为 $m$、电荷量大小为 $q$，粒子重力不计。
\begin{enumerate}
	\item
	求板间电压 $U$ 的大小；
	\item
	若磁场区域如图~\ref{2023河北14b}，磁感应强度不变。调整两板间电压大小，将该粒子从 $P$ 点由静止释放，仍能沿 $\angle CED$ 的角平分线方向从 $E$ 点离开，求粒子从 $C$ 点运动到 $E$ 点的时间。
\end{enumerate}

\twopicture[numA=1, numB=2, labelA=2023河北14a, labelB=2023河北14b, wA=5.5cm, wB=5.5cm, align=right]
{figs/14a.png}
{figs/14b.png}

%% answer:
\jdanswer{
\begin{enumerate}
	\item
	$\frac{qB^{2}a^{2}}{32m}$

	\item
	$\frac{(2\sqrt{3}-3+\pi)m}{3qB}$

\end{enumerate}
}
%% memo:
\memoanswer{
（1）粒子在电场中加速，由动能定理得 $qU=\frac{1}{2}mv^{2}$

根据洛伦兹力提供向心力 $qvB=m\frac{v^{2}}{r}$

根据几何关系有 $r=\frac{a}{4}$

解得板间电压的大小为 $U=\frac{qB^{2}a^{2}}{32m}$

（2）根据题意，做出粒子的运动轨迹，如图所示。

根据几何关系有

$r_{1}=\frac{IK}{\sin 30^\circ}=\frac{\frac{\sqrt{3}}{4}a}{\frac{1}{2}}=\frac{\sqrt{3}}{2}a$

$CK=r_{1}-r_{1}\cos 30^\circ=\frac{\sqrt{3}}{2}a-\frac{3}{4}a$

则

$IJ=IH=\frac{1}{2}FH-CK=\frac{2-\sqrt{3}}{2}a$

粒子在磁场中的运动时间为

$t_{1}=2\times\frac{30^\circ}{360^\circ}\times\frac{2\pi m}{qB}=\frac{\pi m}{3qB}$

粒子在无磁场区域运动的时间为

$t_{2}=\frac{IJ}{v}=\frac{(2\sqrt{3}-3)m}{3qB}$

粒子从 $C$ 点运动到 $E$ 点的时间为

$t=t_{1}+t_{2}=\frac{(2\sqrt{3}-3+\pi)m}{3qB}$

\onepicture[w=4.2cm]{figs/14_an.png}
}

\item
%% number: 15
%% paperName: 2023年河北高考第 15 题 16 分
%% typeId: 6
%% type: 计算题
%% chapter: 运动和力
%% point: 动量和能量
%% method: 无
%% score: 16
%% degree: 350
%% duplicateId: 0
%% body:
如图，质量为 $1\Ukg$ 的薄木板静置于光滑水平地面上，半径为 $0.75\Um$ 的竖直光滑圆弧轨道固定在地面，轨道底端与木板等高，轨道上端点和圆心连线与水平面成 $37^\circ$ 角。质量为 $2\Ukg$ 的小物块 A 以 $8\Ums$ 的初速度从木板左端水平向右滑行，A 与木板间的动摩擦因数为 0.5。当 A 到达木板右端时，木板恰好与轨道底端相碰并被锁定，同时 A 沿圆弧切线方向滑上轨道。待 A 离开轨道后，可随时解除木板锁定，解除锁定时木板的速度与碰撞前瞬间大小相等、方向相反。已知木板长度为 $1.3\Um$，$g$ 取 $10\Umsq$，$\sqrt{10}$ 取 3.16，$\sin 37^\circ=0.6$，$\cos 37^\circ=0.8$。
\begin{enumerate}
	\item
	求木板与轨道底端碰撞前瞬间，物块 A 和木板的速度大小；
	\item
	求物块 A 到达圆弧轨道最高点时受到轨道的弹力大小及离开轨道后距地面的最大高度；
	\item
	物块 A 运动到最大高度时会炸裂成质量比为 $1:3$ 的物块 B 和物块 C，总质量不变，同时系统动能增加 $3\UJ$，其中一块沿原速度方向运动。为保证 B、C 之一落在木板上，求从物块 A 离开轨道到解除木板锁定的时间范围。
\end{enumerate}

\onepicture[w=8.0cm, align=right]{figs/15.png}

%% answer:
\jdanswer{
\begin{enumerate}
	\item
	$v_{1}=7\Ums$，$v_{2}=2\Ums$

	\item
	$F_{N}=\frac{164}{3}\UN$，$H=2\Um$

	\item
	$0.1\Us\leq\Delta t\leq0.118\Us$ 或 $0.732\Us\leq\Delta t\leq0.75\Us$

\end{enumerate}
}
%% memo:
\memoanswer{
（1）设物块 A 的初速度为 $v_{0}$，木板与轨道底部碰撞前，物块 A 和木板的速度分别为 $v_{1}$ 和 $v_{2}$，物块 A、木板的质量分别为 $m_{1}$ 和 $m_{2}$，物块 A 与木板间的动摩擦因数为 $\mu$，木板长度为 $L$，由动量守恒定律和功能关系有

$m_{1}v_{0}=m_{1}v_{1}+m_{2}v_{2}$

$\frac{1}{2}m_{1}v_{0}^{2}=\frac{1}{2}m_{1}v_{1}^{2}+\frac{1}{2}m_{2}v_{2}^{2}+\mu m_{1}gL$

由题意分析 $v_{1}\geqslant v_{2}$，联立式得

$v_{1}=7\Ums$，$v_{2}=2\Ums$

（2）设圆弧轨道半径为 $R$，物块 A 到圆弧轨道最高点时斜抛速度为 $v_{3}$，轨道对物块的弹力为 $F_{N}$。物块 A 从轨道最低点到最高点，根据动能定理有

$-m_{1}gR(1+\sin 37^\circ)=\frac{1}{2}m_{1}v_{3}^{2}-\frac{1}{2}m_{1}v_{1}^{2}$

物块 A 到达圆弧轨道最高点时，根据牛顿第二定律有

$F_{N}+m_{1}g\sin 37^\circ=m_{1}\frac{v_{3}^{2}}{R}$

联立式，得 $F_{N}=\frac{164}{3}\UN$

设物块 A 抛出时速度 $v_{3}$ 的水平和竖直分量分别为 $v_{x}$ 和 $v_{y}$

$v_{x}=v_{3}\sin 37^\circ$，$v_{y}=v_{3}\cos 37^\circ$

斜抛过程物块 A 上升时间 $t_{1}=\frac{v_{y}}{g}=0.4\Us$

该段时间物块 A 向左运动距离为 $s_{1}=v_{x}t_{1}=1.2\Um$

物块 A 距离地面最大高度

$H=R(1+\sin 37^\circ)+\frac{v_{y}^{2}}{2g}=2\Um$

（3）物块 A 从最高点落地时间

$t_{2}=\sqrt{\frac{2H}{g}}=0.632\Us$

设向左为正方向，物块 A 在最高点炸裂为 B、C，设质量和速度分别为 $m_{3}$、$m_{4}$ 和 $v_{4}$、$v_{5}$，设 $m_{3}:m_{4}=1:3$，系统动能增加 $\Delta E_{k}$。根据动量守恒定律和能量守恒定律得

$m_{1}v_{x}=m_{3}v_{4}+m_{4}v_{5}$

$\frac{1}{2}m_{1}v_{x}^{2}+\Delta E_{k}=\frac{1}{2}m_{3}v_{4}^{2}+\frac{1}{2}m_{4}v_{5}^{2}$

解得

$v_{4}=6\Ums$，$v_{5}=2\Ums$ 或 $v_{4}=0$，$v_{5}=4\Ums$

设从物块 A 离开轨道到解除木板锁定的时间范围为 $\Delta t$：

（a）若 $v_{4}=6\Ums$，$v_{5}=2\Ums$，炸裂后 B 落地过程中的水平位移为

$s'=v_{4}t_{2}=3.792\Um$

炸裂后 C 落地过程中的水平位移为

$s''=v_{5}t_{2}=1.264\Um$

木板右端到轨道底端的距离为

$\Delta s=v_{2}(t_{1}+t_{2}-\Delta t)=(2.064-2\Delta t)\Um$

运动轨迹分析如下

\onepicture[w=9.0cm]{figs/15_an1.jpg}

为了保证 B、C 之一落在木板上，需要满足下列条件之一

Ⅰ．若仅 C 落在木板上，应满足 $(2.064-2\Delta t)+1.3\geq1.264+0.6$

且 $(2.064-2\Delta t)\leq1.264+0.6$

解得 $0.1\Us\leq\Delta t\leq0.75\Us$

Ⅱ．若仅 B 落在木板上，应满足 $(2.064-2\Delta t)+1.3\geq3.792+0.6$

且 $(2.064-2\Delta t)\leq3.792+0.6$

不等式无解；

（b）若 $v_{4}=0$，$v_{5}=4\Ums$，炸裂后 B 落地过程中水平位移为 0，炸裂后 C 落地过程中水平位移为

$s'''=v_{5}t_{2}=2.528\Um$

木板右端到轨道底端的距离为

$\Delta s'=v_{2}(t_{1}+t_{2}-\Delta t)=(2.064-2\Delta t)\Um$

运动轨迹分析如下

\onepicture[w=9.0cm]{figs/15_an2.jpg}

为了保证 B、C 之一落在木板上，需要满足下列条件之一

Ⅲ．若仅 B 落在木板上，应满足 $(2.064-2\Delta t)+1.3\geq0.6$

且 $(2.064-2\Delta t)\leq0.6$

解得 $0.732\Us\leqslant\Delta t\leqslant1.382\Us$

Ⅳ．若仅 C 落在木板上，应满足 $(2.064-2\Delta t)+1.3\geq2.528+0.6$

且 $(2.064-2\Delta t)\leq2.528+0.6$

解得 $0\leqslant\Delta t\leqslant0.118\Us$。

综合分析（a）（b）两种情况，为保证 B、C 之一一定落在木板上，$\Delta t$ 满足的条件为 $0.1\Us\leqslant\Delta t\leqslant0.118\Us$ 或 $0.732\Us\leqslant\Delta t\leqslant0.75\Us$
}

\end{enumerate}

\end{document}
